\addcontentsline{toc}{subsection}{1) Formal restatement (\#778)}

Fix $n\in\mathbb{N}$ and let $K_n$ be the complete graph on vertex set $[n]=\{1,\dots,n\}$.  Two players, Alice (red) and Bob (blue), play until every edge of $K_n$ is colored.

\medskip
\noindent\textbf{Game 1 (unbiased clique game $\mathrm{Clique}(n)$).}
Alice moves first. On each move, the player colors one previously uncolored edge with their color.  At the end, let $G_R$ be the red graph and $G_B$ the blue graph.  Alice wins iff
\[
\omega(G_R) > \omega(G_B),
\]
otherwise Bob wins (so a tie counts as a Bob win).

\medskip
\noindent\textbf{Question 1.} Does Bob have a winning strategy for every $n\ge 3$?

\medskip
\noindent\textbf{Game 2 (biased clique game $\mathrm{Clique}^{(1,2)}(n)$).}
Alice colors one new edge; then Bob colors two new edges; repeat until all edges are colored.
At the end Alice wins iff $\omega(G_R) \ge \omega(G_B)$ (equivalently Bob wins iff $\omega(G_B) > \omega(G_R)$).

\medskip
\noindent\textbf{Question 2.} For $n>3$, does Bob have a winning strategy?

\medskip
\noindent\textbf{Game 3 (maximum-degree / ``star'' game $\mathrm{Star}(n)$).}
This is the unbiased game as in Game~1 (one edge per move, Alice first), but Alice wins iff
\[
\Delta(G_R) > \Delta(G_B),
\]
otherwise Bob wins.

\medskip
\noindent\textbf{Question 3.} Who wins $\mathrm{Star}(n)$?

\medskip
\noindent\textbf{Useful reformulations.}
Because every edge is colored exactly once, $G_B$ is the complement of $G_R$ in $K_n$.
Hence
\[
\omega(G_B)=\omega(\overline{G_R})=\alpha(G_R),
\]
so in Game~1 Alice wins iff $\omega(G_R) > \alpha(G_R)$.
Similarly in Game~3,
\[
\Delta(G_B)=\Delta(\overline{G_R}) = (n-1) - \delta(G_R),
\]
so Alice wins iff
\[
\Delta(G_R) + \delta(G_R) > n-1.
\]
The game is still genuinely two-player: Bob does not control $G_R$ directly, but he controls which edges remain available.
